% Programming with InitechBase, Chapter 5: Functions for Fields and Memory
% Variables (AN-16; annex Sec 8.2 P5).  Panels replayed (examples/p5-*.dot);
% sources core/fn_builtins.c (STR, LTRIM, TRANSFORM, DTOC, CTOD, CMONTH,
% IIF in core/eval.c).
\chapter{Functions for Fields and Memory Variables}
\Gates{test-pubs-examples, test-xbase-fn-a, test-xbase-fn-c, test-xbase-transform}
A program prints values in columns, joins them into lines, and converts them
from one type to another.  Functions do this work.  This chapter shows the
functions the Initech Accounting System uses to build its lines; Chapter~6 of
\doctitle{Using InitechBase} describes every function.

\begin{ChapterCovers}
\item Joining values into a line
\item Numbers as characters: \fn{STR()}, \fn{LTRIM()}, \fn{TRANSFORM()}
\item Dates as characters: \fn{DTOC()} and \fn{CTOD()}
\item Choosing a value: \fn{IIF()}
\item A formatted register
\end{ChapterCovers}

\ATEHead{Preparing for This Chapter}
Start InitechBase with a fresh copy of the sample diskette, and type
\type{USE INVOICE}.

\ATEHead{Joining Values}
\texttt{+} joins character values into one, which \cmd{?} then prints as a single
line.  Every part must be a character value: InitechBase does not join a number
or a date to characters.  Fields of type Character are joined as they are,
blanks and all, which keeps the columns of a listing in line.

\ATEHead{Numbers as Characters}
\fn{STR()} converts a number to characters of a fixed width, with a fixed number
of decimal places, right-aligned; it is the function that lines numbers up in
a column.  \fn{LTRIM()} removes the leading blanks when the number is to follow
a word.  \fn{TRANSFORM()} converts a number by a picture, in which each 9 holds a
digit:
\DotExample{examples/p5-numbers}

\ATEHead{Dates as Characters}
\fn{DTOC()} converts a date to characters in the form set by SET DATE;
\fn{CTOD()} converts characters to a date.  A date plus a number of days is a
date; \fn{CMONTH()} gives the name of its month:
\DotExample{examples/p5-dates}
The year of a date has two digits unless SET CENTURY is ON.  Two digits typed
into \fn{CTOD()} are a year of the nineteen hundreds; this is the date
convention of every program of the system.

\ATEHead{Choosing a Value}
\fn{IIF()} returns its second argument when its first is true, and its third when
it is false.  It chooses a word for a logical field without an \cmd{IF}.

\ATEHead{A Formatted Register}
\file{INVFMT.PRG} prints the ledger in fixed columns, with the status of each
invoice in words:
\ProgramListing{INVFMT.PRG}
\DotExample{examples/p5-invfmt}
Each line is built from the fields with \texttt{+}, in two steps to keep the
program lines short: the memory variable \texttt{MLINE} holds the first part.
\fn{DTOC()} converts the due date, \fn{TRANSFORM()} the amount, and \fn{IIF()}
turns \texttt{PAID} into a word, kept in the memory variable \texttt{MSTAT}.

\ATEHead{Where Do You Go From Here?}
Chapter~6 lets the program ask the user for what it needs.

\ATEHead{Summary}
\begin{enumerate}
\item \texttt{+} joins character values; numbers and dates are converted first.
\item \fn{STR()}, \fn{LTRIM()} and \fn{TRANSFORM()} turn numbers into characters.
\item \fn{DTOC()} and \fn{CTOD()} convert between dates and characters.
\item \fn{IIF()} chooses between two values.
\end{enumerate}
